\documentclass[a4paper,11pt,twoside]{article}

%-----------------------------------
%
%    CERMICS 2009
%  Last modified Gabriel STOLTZ
%
%----------------------------------

\usepackage[latin1]{inputenc} %Ligne PS
\usepackage[T1]{fontenc} %Ligne PS
\newlength{\minipagewidth}
\setlength{\minipagewidth}{\textwidth}
\setlength{\fboxsep}{3mm}
\usepackage{a4wide}
%\usepackage[active]{srcltx}
\usepackage{amsmath,amsthm}
\usepackage{amssymb}
\usepackage{a4wide}
\usepackage{graphicx}
\usepackage{color}
\usepackage{epic}
\usepackage{enumerate}


\bibliographystyle{plain}
\newcommand{\dps}{\displaystyle}

%-------------
%   d�but
%-------------

\begin{document}

%----------------
%  TITRE
%----------------

\title{Elements on the ``dimer in a WCA solvent'' code}
\author{Gabriel Stoltz ({\tt stoltz@cermics.enpc.fr})}

\maketitle

\section{Overall structure}

The source files are in the directoty {\tt src/}, the input files are in the directory
{\tt input/}, the output files are in the directory {\tt output/}, and the present
documentation is in {\tt doc/}. The codes should be compiled and run within the 
{\tt src/} directory as they stand. The input file read by the executable 
should be self-explanatory, otherwise the user should have a look at the corresponding
file {\tt input$\_$METHOD.cpp} which has been compiled depending on the method
(see the {\tt makefile} command file for a list of files required for the compilation).

Besides header files (extension {\tt .hpp}), the source files
include {\tt gaussian.cpp} which is used to produce gaussian random numbers; 
{\tt vector.cpp} and {\tt matrix.cpp} which are simple vector and matrix classes
defined in order to require the least amount of {\tt C++} libraries; {\tt main.cpp} which 
is the (concise) main files, in which a call to the function {\tt load} is performed.
This is the master function in the codes, which is defined in a file whose name
depends on the method at hand; finally, forces and energy computations are
performed with {\tt hamiltonian.cpp}.

\section{Model and input files}

\subsection{Microscopic interactions and thermodynamic conditions}

The parameters describing the microscopic interactions and the thermodynamic
state of the system are the same for all input files:
\begin{itemize}
\item {\tt Friction coefficient} $\gamma$ used in the Langevin dynamics;
\item {\tt Inverse temperature} $\beta$;
\item {\tt Number of particles} $n$, the total number being then $N=n^2$, 
  and the number of solvent particles is $N-2$;
\item {\tt Elementary cell size} $a$, so that the density of solvent is 
  $\rho = (1-2/N) a^{-2}$;
\item {\tt WCA equilibrium distance} $\sigma$;
\item {\tt WCA energy epsilon} $\varepsilon$;
\item {\tt Dimer barrier height} $h$;
\item {\tt Dimer bond length} $w$.
\end{itemize}
Besides, time step $\Delta t$ (see the field {\tt Time step}) is needed
in all input files.

\subsection{Canonical sampling}

A simple canonical sampling can be obtained by {\tt make sampling}, which produces
{\tt sampling++}. The output files are: {\tt energy} gives the energy 
as a function of time; {\tt coord} is the instantaneous distribution of 
the values of the reaction coordinate; {\tt xmakemol.xyz} is a movie of the system
readable with the software {\tt XMakeMol} (see {\tt www.nongnu.org/xmakemol/}).

The additional parameters are
\begin{itemize}
\item {\tt Dynamics (0=Ovrdmpd, 1=Lngvn)} depending on whether Langevin dynamics
or overdamped dynamics should be used; 
\item {\tt Number of iterations};
\item {\tt Xmakemol output frequency} is the number of steps before a new 
configurations is written down;
\item {\tt Other outputs frequency} is the frequency at which the energy and the value
of the reaction coordinate are saved.
\end{itemize}

\subsection{Extended bridge sampling}

The code is compiled with the instruction {\tt make mbar}, which produces
{\tt mbar++}.
The additional parameters for the production part are
\begin{itemize}
\item {\tt Number config. / restraint} is the required number of configurations per 
restraining center;
\item {\tt Preliminary thermalization} is the number of steps of the dynamics before 
configurations are acquired; 
\item {\tt Number of subsampling steps} is the number of steps before a new configuration
is saved;
\item {\tt Restraining potential} $K$ is spring constant for the restraining potential; 
\item {\tt Lower restraint center} $z_{\rm min}$ is the opsition of the lowest center; 
\item {\tt Upper restraint center} $z_{\rm max}$ is the position of the largest center;
\item {\tt Number of restraint cent.} gives the number $N_z$ of centers considered;
\end{itemize}
while the parameters for the post-processing part are
\begin{itemize}
\item {\tt Lower value of RC} $\xi_{\rm min}$ and \item {\tt Upper value of RC} $\xi_{\rm max}$
give the range of values of the reaction coordinate for which the free energy is required; 
\item {\tt Number of PMF RC values} $N_\xi$ gives the number of bins in this interval; 
\item {\tt Tolerance for SC convergence} is the stopping criterion for the
  simple fixed-point iteration used to solve the nonlinear equation
  on the ratios of normalizing constants.
\end{itemize}
There are several output files: {\tt energy} and {\tt coord} list, in each line, the
energies and the values of the reaction coordinate for the configurations sampled
at a given value of the restraining center; {\tt ratio} gives the final ratios
of partition functions ($N_z$ values), while {\tt PMF} is the final potential
of mean force profile in the interval $[x_{\rm min}, \xi_{\rm max}]$, obtained
by computing the average of indicator functions of bins of size
$\Delta z = (\xi_{\rm max}-\xi_{\rm min})/N_\xi$.

\subsection{Thermodynamic integration}

The code is compiled with the instruction {\tt make ti}, which produces {\tt ti++}.
The parameters of the dynamics are
\begin{itemize}
\item {\tt Dynamics (0=Ovrdmpd, 1=Lngvn)} depending on whether projected Langevin dynamics
or projected overdamped dynamics should be used;
\item {\tt Number of iterations};
\item {\tt Number of thermalization steps}: number of steps performed
  each time the value of the reaction coordinate changes before the computation of the 
  approximation of the mean force by ergodic averages starts; 
\item {\tt Output frequency} is the number of time steps before the outputs are written out
  (see below for a list);
\item {\tt Lower value of RC} and {\tt Upper value of RC} define the biasing region where
the mean force is estimated (the bias being 0 outside this zone);  
\item {\tt Number of PMF RC values} gives the number of bins in which the reaction coordinate
space is separated; 
\end{itemize}
The output files are the following:
\begin{itemize}
\item {\tt coord} lists the values of the reaction coordinate;
\item {\tt energy} gives the total energy and the potential energy;
\item {\tt current$\_$mean$\_$force} is the current ergodic approximation of the mean force.
In the overdamped case, the first column is 
obtained by averaging the local mean force (known analytically here), the second
by averaging the plain Lagrange multipliers, the third by averaging the Lagrange multipliers 
with elimination of the martingale part, the fourth with the variance reduction method
using time reversed increments. In the Langevin case, 
the first column is obtained by averaging the local mean force, the second one by averaging
the Lagrange multipliers used in the Rattle part, the third by computing the momentum
average local mean force, the fourth and fifth ones being
the averages of the Lagrange multipliers of the position and momentum constraints, 
respectively.
\item {\tt Lagrange} gives the instantaneous values of the mean force and Lagrange multipliers
which are averaged in the previous file (same ordering).
\item {\tt mean$\_$force} and {\tt PMF} give the final mean force and free energy profiles
respectively.
\end{itemize}

\subsection{Adaptive dynamics}

The code is compiled with the instruction {\tt make abf}, which produces
{\tt abf++}. There are several output files: {\tt energy} gives the energy 
as a function of time; {\tt coord} is the instantaneous distribution of 
the values of the reaction coordinate, while {\tt coord$\_$first} is the same information
for the first system only; {\tt mean$\_$force} is the current mean force (the values
of the reaction coordinate bins are in the first line, and each subsequent line
is a profile some time later than the previous one), while {\tt PMF} is the current
associated free energy profile obtained by numerical integration; finally,
{\tt histo} is the cumulated histogram of visits in a bin.

The parameters of the dynamics are
\begin{itemize}
\item {\tt Dynamics (0=Ovrdmpd, 1=Lngvn)} depending on whether Langevin dynamics
or overdamped dynamics should be used;
\item {\tt Number of iterations};
\item {\tt Number of replicas} gives the number of systems simulated in parallel 
but contributing to the same free energy profile;
\item {\tt Mean force/bias output freq.} gives the frequency at which mean force,
free energy profiles, distribution of reaction coordinates and cumulated histogram
are written out (counted in number of iteration steps)
\item {\tt Other outputs frequency} is the number of time steps before the current 
energy and reaction coordinate value of the first system are written out;
\item {\tt Lower value of RC} and {\tt Upper value of RC} define the biasing region where
the mean force is estimated (the bias being 0 outside this zone);  
\item {\tt Number of PMF RC values} gives the number of bins in which the reaction coordinate
space is separated; 
\item {\tt Selection} is set to 1 when selection is required; 
\item {\tt Section intensity} is the prefactor before the selection term (provided selection
is required). 
\end{itemize}

\subsection{Nonequilibrium dynamics}

The code is compiled with the instruction {\tt make jarz}, which produces
{\tt jarz++}.

\begin{itemize}
\item {\tt Number of thermalization steps } is the subsampling time used to sample
the initial conditions from a projected dynamics;
\item {\tt Work distribution output freq.} is the number of steps 
after which the whole work distribution is written out;
\item {\tt Other output frequency} is the number of steps after which the energies
and the value of the reaction coordinate are saved;         
\item {\tt Lower value of RC} is the initial value of the reaction coordinate
at the beginning of the switching;     
\item {\tt Upper value of RC} is the final value of the reaction coordinate;     
\item {\tt Switching time} gives the time to perform the switching;            
\item {\tt Number of replicas}           
\end{itemize}
The output files are the following:
\begin{itemize}
\item {\tt coord} gives the value of the reaction coordinate of the first system
as a function of time;
\item {\tt energy} gives the total energy and the potential energies as a function of time; 
\item {\tt PMF} gives the free energy profile as a function of time, with the estimate
using the local mean force on the second column, and the estimates obtained with the 
Lagrange multipliers on the third one;
\item {\tt Work} is the work distribution at the required times, for the works
computed with the local mean force;      
\item {\tt Work$\_$lagrange} is the work distributions obtained for the works computed
from the Lagrange multipliers;     
\item {\tt current$\_$work} gives the current work for the first system, computed using the 
local mean force on the first column, and the Lagrange multipliers on the second.
\end{itemize}

\section{Useful analytical expressions}

\paragraph{Reaction coordinate and derivatives.}
The chosen reaction coordinate is
\[
\xi(q) = \frac{|q_1-q_2|-r_0}{2w},
\]
where $q_1,q_2 \in \mathrm{R}^d$, 
$d$ denoting the dimension of the ambient physical space ($d=2$ here).
Its gradient is
\[
\nabla \xi(q) = \frac{1}{2w}\left( \begin{array}{c}
\dps \frac{q_1-q_2}{|q_1-q_2|} \\ [10pt]
\dps -\frac{q_1-q_2}{|q_1-q_2|} \\
0 \\
\vdots \\
0
\end{array} \right),
\]
and 
\[
|\nabla \xi(q)|^2 = \frac{1}{2w^2}.
\]
The local mean force is
\begin{eqnarray*}
f(q) & = & \frac{\nabla \xi(q) \cdot \nabla V(q)}{|\nabla \xi(q)|^2} - \frac1\beta \,  
{\rm div}\left( \frac{\nabla \xi(q)}{|\nabla \xi(q)|^2} \right) \\
& = & \frac{w}{|q_1-q_2|} 
\left [ (q_1-q_2) \cdot \Big(\partial_{q_1} V(q) - \partial_{q_2} V(q) \Big)
- \frac{2(d-1)}{\beta} \right ].\\
\end{eqnarray*}
Finally, notice that there is no Fixman term since $|\nabla \xi|$ is constant.

\paragraph{Constrained overdamped Langevin processes.}
The constrained dynamics
\[
\left \{ \begin{array}{rcl}
dq_t & = & \dps -\nabla V(q_t) \, dt + \sqrt{\frac{2}{\beta}} \, dW_t 
+ \nabla \xi(q_t) \, d\Lambda_t, \\ [5pt]
\xi(q_t) & = & z,
\end{array} \right.
\]
is integrated numerically with the scheme
\[
\left \{ \begin{array}{rcl}
q^{n+1} & = & \dps q^n -\nabla V(q^n) \, \Delta t + \sqrt{\frac{2 \Delta t}{\beta}} \, G^n 
+ \nabla \xi(q^{n+1}) \, \Delta \Lambda^{n+1}, \\ [5pt]
\xi(q^{n+1}) & = & z.
\end{array} \right.
\]
Therefore, 
\[
\left( 1 - \frac{\Delta \Lambda^{n+1}}{w|q_2^{n+1}-q_1^{n+1}|}\right) 
(q_2^{n+1}-q_1^{n+1}) = \widetilde{q}_2^n - \widetilde{q}_1^n,
\]
with 
\[
\widetilde{q}^n = q^n -\nabla V(q^n) \, \Delta t + \sqrt{\frac{2 \Delta t}{\beta}} \, G^n.
\]
and $G^n = (G_1^n,G_2^n,\dots)$ with $G_i^n \in \mathbb{R}^d$.
Since $|q_2^{n+1}-q_1^{n+1}| = 2wz + r_0$ is fixed,
the lagrange multiplier reads
\[
\Delta \Lambda^{n+1} = w\Big ( 2wz + r_0 - \left|\widetilde{q}_2^n - \widetilde{q}_1^n\right|\Big )
= 2w^2\Big (z-\xi(\widetilde{q})\Big ).
\]
Besides, $q_1^{n+1}+q_2^{n+1} = \widetilde{q}_2^n + \widetilde{q}_1^n$, so that
\[
q_1^{n+1} = \frac12 \left( 1 + \frac{1}{1-\Delta \Lambda^{n+1} / [w(2wz+r_0)]}\right ) 
\, \widetilde{q}_1^n
+ \frac12 \left( 1 - \frac{1}{1-\Delta \Lambda^{n+1} / [w(2wz+r_0)]}\right ) 
\, \widetilde{q}_2^n,
\]
and a similar expression for $q_2^{n+1}$. Some care is however required to treat
correctly the periodic boundary conditions.
With the variance reduction procedure, 
\[
\Delta \widetilde{\Lambda}^{n+1} = \Delta \Lambda^{n+1} + \sqrt{\frac{2\Delta t}{\beta}} \,
\frac{\nabla \xi(q^n) \cdot G^n}{|\nabla \xi(q^n)|^2}
=  \Delta \Lambda^{n+1} + w \frac{q_1^n-q_2^n}{|q_1^n-q_2^n|} \cdot (G_1^n - G_2^n).
\]
The expressions in the case of nonequilibrium switching processes are similar.
It suffices to consider an additional term related to the variation of the constraint:
\[
\left \{ \begin{array}{rcl}
q^{n+1} & = & \dps q^n -\nabla V(q^n) \, \Delta t + \sqrt{\frac{2 \Delta t}{\beta}} \, G^n 
+ \nabla \xi(q^{n+1}) \, \Delta \Lambda^{n+1}, \\ [5pt]
\xi(q^{n+1}) & = & z^{n+1},
\end{array} \right.
\]
and 
\[
\Delta \Lambda_{\rm noneq}^{n+1} = \Delta \Lambda^{n+1} + \sqrt{\frac{2\Delta t}{\beta}} \,
\frac{\nabla \xi(q^n) \cdot G^n}{|\nabla \xi(q^n)|^2}
=  \Delta \Lambda^{n+1} + w \frac{q_1^n-q_2^n}{|q_1^n-q_2^n|} \cdot (G_1^n - G_2^n)
- \frac{z^{n+1}-z^n}{|\nabla \xi(q^n)|^2}.
\]

\paragraph{Constrained Langevin processes.}
We detail how the Rattle step 
\[
\left \{ \begin{array}{rclc}
p^{n+1/2} & = & \dps p^{n} - \frac{\Delta t}{2} \nabla V (q^{n}) 
+ \nabla \xi(q^n) \, \lambda^{n+1/2}, & \\ [5pt]
\xi(q^{n+1}) & = & z, & (C_{q})\\[5pt]
q^{n+1} & = & q^{n} + \Delta t \, M^{-1} \, p^{n+1/2}, & \\[5pt]
p^{n+1} & = & \dps p^{n+1/2} - \frac{\Delta t}{2} \nabla V (q^{n+1}) 
+ \nabla \xi(q^{n+1}) \, \lambda^{n+1}, &\\[5pt]
\nabla\xi (q^{n+1})^T M^{-1} p^{n+1} & = & 0, & (C_{p})
\end{array} \right.
\]
is implemented. In the code, $M = \mathrm{Id}$. Denoting 
\[
\widetilde{q}^n = q^{n} + \Delta t \, p^n -\frac{\Delta t^2}{2} \, \nabla V(q^n),
\] 
the value of $\lambda^{n+1/2}$ is determined by the condition 
\[
\xi\left( \widetilde{q}^n + \Delta t \, \lambda^{n+1/2} \nabla \xi(q^n) \right) = z,
\]
which can be rewritten as
\[
\left\| \widetilde{q}_1^n - \widetilde{q}_2^n + \frac{\Delta t \, \lambda^n}{w} \, 
e_{12}^n \right\|^2 = z^2,
\]
with 
\[
e_{12}^n \equiv e_{12}(q^n) =  \frac{q_1^n-q_2^n}{|q_1^n-q_2^n|}.
\]
This equation is a second-order equation in the $\Delta t \, \lambda$ variable:
\[
(\Delta t \,\lambda)^2 + 2 w e_{12}^n \cdot \left(\widetilde{q}_1^n - \widetilde{q}_2^n\right) 
\Delta t \, \lambda
+ w^2 \left( \left \| \widetilde{q}_1^n - \widetilde{q}_2^n \right \|^2 - z^2\right) = 0,
\]
whose solutions are
\[
\lambda_\pm = \frac{w}{\Delta t} 
\Big ( -e_{12}^n \cdot \left(\widetilde{q}_1^n - \widetilde{q}_2^n\right)
\pm \sqrt{\Delta} \Big ),
\]
with 
\[
\Delta = \Big [ e_{12}^n \cdot \left(\widetilde{q}_1^n - \widetilde{q}_2^n\right) \Big ]^2
- 4 \left( \left \| \widetilde{q}_1^n - \widetilde{q}_2^n \right \|^2 - z^2\right).
\]
The smallest $\lambda$ is chosen, so that $\lambda^n = \lambda_-$ when 
$e_{12}^n \cdot \left(\widetilde{q}_1^n - \widetilde{q}_2^n\right) \geq 0$ and
$\lambda^n = \lambda_+$ otherwise.

The Lagrange multiplier $\lambda^{n+1/2}$ is determined by 
\[
\nabla \xi(q^{n+1})^T M^{-1} p^{n+1} = 0,
\]
which can be rewritten as
\[
e_{12}^{n+1} \cdot \left(p_1^{n+1} - p_2^{n+1}\right) =
e_{12}^{n+1} \cdot \left(\widetilde{p}_1^{n+1} - \widetilde{p}_2^{n+1} + 
\frac{\lambda^{n+1}}{w} e_{12}^{n+1} \right) = 0,
\]
with 
\[
\widetilde{p}^{n+1} = p^{n+1/2} - \frac{\Delta t}{2} \nabla V(q^{n+1}).
\]
Then,
\[
\lambda^{n+1} = -w \, e_{12}^{n+1} \cdot \left(\widetilde{p}_1^{n+1} - 
\widetilde{p}_2^{n+1}\right).
\]

The comparison with the local rigid mean force (equal to the local mean force)
requires the computation of 
\[
f_{\rm rgd}(q,p) = G_M^{-1} \nabla \xi^T M^{-1}\nabla V(q) 
- G_M^{-1}(q) {\rm Hess}_q(\xi)(M^{-1} p, M^{-1} p).
\]
The first term is the same as for the local mean force. For the second term,
some straightforward computations give
\begin{eqnarray*}
& & G_M^{-1}(q) {\rm Hess}_q(\xi)(M^{-1} p, M^{-1} p) 
= w \nabla^2 \xi(q) : p \otimes p \\
& & = \frac{w}{|q_1-q_2|} \left[ (p_{1,x}-p_{2,x})^2 + (p_{1,y}-p_{2,y})^2 - 
\Big(\frac{x_1-x_2}{|q_1-q_2|} \, (p_{1,x}-p_{2,x})
+\frac{y_1-y_2}{|q_1-q_2|} \, (p_{1,y}-p_{2,y})\Big)^2\right].
\end{eqnarray*}
Notice that
\[
\frac{x_1-x_2}{|q_1-q_2|} \, (p_{1,x}-p_{2,x})
+\frac{y_1-y_2}{|q_1-q_2|} \, (p_{1,y}-p_{2,y})
= 2w \, e_{12} \cdot (p_1-p_2).
\]
This quantity is therefore equal to 0 for constrained processes
(and in general proportional to the momentum or velocity constraint).


\paragraph{Nonequilibrium Langevin processes.}
We describe how the momentum constrained is handled. It reads, in the nonequilibrium case,
\[
\nabla\xi (q^{n+1})^T M^{-1} p^{n+1} = \frac{z(t_{n+1})-z(t_n)}{\Delta t}.
\]
with 
\[
p^{n+1} = p^{n+1/2} - \frac{\Delta t}{2} \nabla V (q^{n+1}) 
+ \nabla \xi(q^{n+1}) \, \lambda^{n+1}.
\]
Then,
\[
\lambda^{n+1} = w \left[ -e_{12}^{n+1} \cdot \left(\widetilde{p}_1^{n+1} - 
\widetilde{p}_2^{n+1}\right) + \frac{z(t_{n+1})-z(t_n)}{\Delta t} \right]
\]
with 
\[
\widetilde{p}^{n+1} = p^{n+1/2} - \frac{\Delta t}{2} \nabla V(q^{n+1}).
\]

\end{document}
